\slajdid{09_3-s070}
\begin{frame}[label=09_3-s070]{Domino logika}
\gusttekst\setlength{\abovedisplayskip}{3pt}\setlength{\belowdisplayskip}{3pt}\setlength{\abovedisplayshortskip}{2pt}\setlength{\belowdisplayshortskip}{2pt}Kod dinamičke logike ozbiljan problem leži prilikom kaskadne veze funkcija. Da vidimo to na jednostavnom primeru
\par\medskip\begin{columns}[c,onlytextwidth]\column{.4\textwidth}\centering\begin{tikzpicture}[x=.6cm,y=.55cm,font=\sitneoznake,>=Latex]\ctikzset{bipoles/length=.6cm}\begin{scope}[shift={(0,3)}]\draw(0,-.35)--(0,.35);\draw(-.13,-.4)--(-.13,.4);\draw(0,.3)--(.4,.3)--(.4,.65);\draw(0,-.3)--(.4,-.3)--(.4,-.65);\draw(-.23,0)circle(.1);\draw(-1,0)node[left]{}--(-.33,0);\node[right]at(.52,0){P1};\end{scope}\begin{scope}[shift={(0,0.6)}]\draw(0,-.35)--(0,.35);\draw(-.13,-.4)--(-.13,.4);\draw(0,.3)--(.4,.3)--(.4,.65);\draw(0,-.3)--(.4,-.3)--(.4,-.65);\draw(-1,0)node[left]{}--(-.13,0);\node[right]at(.7,-.35){T1};\end{scope}\begin{scope}[shift={(0,-1)}]\draw(0,-.35)--(0,.35);\draw(-.13,-.4)--(-.13,.4);\draw(0,.3)--(.4,.3)--(.4,.65);\draw(0,-.3)--(.4,-.3)--(.4,-.65);\draw(-1,0)node[left]{}--(-.13,0);\node[right]at(.52,0){N1};\end{scope}\draw(0.4,3.65)--(0.4,4.1);\draw(0.1,4.1)--(0.7,4.1);\node[above]at(0.4,4.1){$V_{DD}$};\draw(0.4,2.35)--(0.4,1.25);\draw(0.4,-.05)--(0.4,-.35);\draw(0.4,-1.65)--(0.4,-2)node[ground]{};\begin{scope}[shift={(3.5,3)}]\draw(0,-.35)--(0,.35);\draw(-.13,-.4)--(-.13,.4);\draw(0,.3)--(.4,.3)--(.4,.65);\draw(0,-.3)--(.4,-.3)--(.4,-.65);\draw(-.23,0)circle(.1);\draw(-1,0)node[left]{}--(-.33,0);\node[right]at(.52,0){P2};\end{scope}\begin{scope}[shift={(3.5,0.6)}]\draw(0,-.35)--(0,.35);\draw(-.13,-.4)--(-.13,.4);\draw(0,.3)--(.4,.3)--(.4,.65);\draw(0,-.3)--(.4,-.3)--(.4,-.65);\draw(-1,0)node[left]{}--(-.13,0);\node[right]at(.7,-.35){T2};\end{scope}\begin{scope}[shift={(3.5,-1)}]\draw(0,-.35)--(0,.35);\draw(-.13,-.4)--(-.13,.4);\draw(0,.3)--(.4,.3)--(.4,.65);\draw(0,-.3)--(.4,-.3)--(.4,-.65);\draw(-1,0)node[left]{}--(-.13,0);\node[right]at(.52,0){N2};\end{scope}\draw(3.9,3.65)--(3.9,4.1);\draw(3.6,4.1)--(4.2,4.1);\node[above]at(3.9,4.1){$V_{DD}$};\draw(3.9,2.35)--(3.9,1.25);\draw(3.9,-.05)--(3.9,-.35);\draw(3.9,-1.65)--(3.9,-2)node[ground]{};\draw(-1,3)--(-1.7,3)--(-1.7,-1)--(-1,-1);\draw(-1.7,-1)--(-2.2,-1)node[ocirc]{}node[above]{CLK};\draw(-1,.6)--(-1.8,.6)node[ocirc]{}node[left]{A};\draw(-1.3,-1)--(-1.3,-2.6)--(1.8,-2.6)--(1.8,3)--(2.5,3);\draw(1.8,-1)--(2.5,-1);
\draw(.4,1.9)--(1.3,1.9)--(1.3,.6)--(1.65,.6);\draw[preaction={draw=white,line width=2pt}](1.65,.6)arc[start angle=180,end angle=0,x radius=.15,y radius=.12];\draw(1.95,.6)--(2.5,.6);\draw(3.9,1.9)--(6.1,1.9)node[ocirc]{}node[right]{Y};\draw(5.5,1.9)to[C,l=$C_L$](5.5,.1)node[ground]{};
\end{tikzpicture}
\column{.58\textwidth}
Ako je A na logičkoj nuli, u fazi pripreme kapacitivnost tranzistora T2 će se napuniti na $V_{DD}$. U fazi izračunavanja tranzistor T1 je zakočen, tranzistor T2 će vodi i prazniće izlaznu kapacitivnost, tako da će na kraju faze izračunavanja biti logička nula na izlazu kao što i treba da bude. Dvostruka inverzija.
\par\medskip
Ako je A na logičkoj jedinici, u fazi pripreme kapacitivnost tranzistora tranzistora T2 će se napuniti na $V_{DD}$. U fazi izračunavanja tranzistor T1 je voditi i prazniće kapacitivnost tranzistora T2. Međutim tranzistor T2 će voditi dok god mu se ne isprazne njegove kapacitivnosti a samim time će prazniti i izlaznu kapacitivnost. Može da se desi da na kraju faze izračunavanja napon na izlazu padne ispod logičke jedinice koja bi trebala da bude na izlazu.
\end{columns}
\end{frame}
